\subsection{Target Users and Environments}

The system mainly supports two user groups:

\begin{itemize}
    \item \textbf{Patients or general users:}
    People who want to manage their medications more safely, check possible medication-related risks, and follow their medication schedules more reliably.

    \item \textbf{Healthcare professionals:}
    Doctors, pharmacists, nurses, or other qualified professionals who can review high-risk or uncertain medication-related cases and provide professional guidance.
\end{itemize}

The system is a responsive Web application for desktop and mobile browsers.

\subsection{Scope and Constraints}

The project focuses on \textbf{helping users understand medication risks and manage confirmed medication schedules}.

The core scope includes:

\begin{itemize}
    \item Extracting prescription information using OCR, supporting manual entry, and identifying user requests through a chatbot or feature selection.

    \item Supporting user verification of extracted details and identifying medications and active ingredients using reference data.

    \item Automatically checking drug-drug interactions and severity using DDInter 2.0; therapeutic duplication and drug-disease checks depend on data availability.

    \item Presenting safety findings, limitations, and guidance, with professional review for high-risk or uncertain cases.

    \item Creating reminders from confirmed prescription schedules, with basic adherence tracking and a safety gate for severe interactions.

    \item Recording possible adverse drug reactions for follow-up and professional review when necessary.
\end{itemize}

The main constraints are:

\begin{itemize}
    \item The system does not independently prescribe, stop, replace, or modify prescription medications.

    \item Safety assessment requires user-verified medication details and reliable identification. Unresolved or insufficiently assessed cases must be reported as unassessed, not safe.

    \item Assessment coverage and permitted use depend on the reliability and licensing of available medication data sources. No detected interaction does not guarantee safety.

    \item New reminders for medications involved in a detected severe interaction require professional review and confirmation; unrelated medications may proceed. Existing treatment or reminders are not automatically stopped.
    
    \item Medication reminders follow user-confirmed prescription instructions. Changes to prescribed dosing schedules require healthcare professional confirmation and user acceptance.

    \item Reported symptoms indicate possible adverse drug reactions, not confirmed causation.

    \item The prototype supports printed or digital Vietnamese prescriptions and simulates professional review without integration with real clinical services.
\end{itemize}

Advanced timing analysis, drug-food checks, pill recognition, and handwritten prescription support are outside the initial scope.
